\documentclass[a4paper,11pt]{article}
\usepackage[utf8]{inputenc}
\usepackage[T1]{fontenc}
\usepackage[french]{babel}
\usepackage{amsmath}
\usepackage[top=2.2cm, bottom=2.2cm, left=2.3cm, right=2.3cm]{geometry}
\usepackage{xcolor}
\usepackage{booktabs}
\usepackage{tabularx}
\usepackage{enumitem}
\usepackage{listings}
\usepackage{titlesec}
\usepackage[most]{tcolorbox}
\usepackage{hyperref}
\hypersetup{colorlinks=true, linkcolor=primary, urlcolor=primary}

\definecolor{primary}{HTML}{1A237E}
\definecolor{success}{HTML}{2E7D32}
\definecolor{warning}{HTML}{E65100}
\definecolor{codebg}{HTML}{F6F6F6}

\lstset{backgroundcolor=\color{codebg}, basicstyle=\ttfamily\footnotesize,
  breaklines=true, frame=single, framesep=4pt, showstringspaces=false,
  columns=fullflexible, keywordstyle=\color{primary}}

\titleformat{\section}{\large\bfseries\color{primary}}{\thesection.}{0.6em}{}
\titlespacing*{\section}{0pt}{13pt}{5pt}
\titleformat{\subsection}{\bfseries}{\thesubsection.}{0.5em}{}
\titlespacing*{\subsection}{0pt}{8pt}{3pt}
\setlength{\parindent}{0pt}
\setlength{\parskip}{4pt}
\newcommand{\code}[1]{\texttt{#1}}

\newtcolorbox{note}[1][Note]{colback=blue!4, colframe=primary, fonttitle=\bfseries,
  title=#1, boxrule=0.5pt, left=6pt, right=6pt, top=4pt, bottom=4pt}

\title{\vspace{-1cm}\textbf{Communication V2V coopérative : CAM et DENM\\[3pt]
\large Se signaler, percevoir, alerter --- de la position au freinage d'urgence}}
\author{Douae Sebti}
\date{10 juillet 2026}

\begin{document}
\maketitle
\vspace{-6pt}\hrule\vspace{8pt}

\begin{note}[Vue d'ensemble]
Deux véhicules communiquent par messages ETSI sur WiFi. Le \textbf{CAM} dit en continu
« je suis là » (position, cap, vitesse). Le \textbf{DENM} dit, sur événement,
« attention, danger ici ». En reliant une \textbf{perception} (caméra + IA) à l'émission
d'un DENM, j'obtiens une boucle de sécurité coopérative complète : \emph{percevoir →
alerter → réagir}.
\end{note}

\tableofcontents
\vspace{6pt}\hrule

% =====================================================================
\section{Le système V2V coopératif}
Un système \textbf{V2V} (\emph{vehicle-to-vehicle}) permet aux véhicules de se parler pour
coopérer. La norme européenne \textbf{ETSI C-ITS} définit deux messages complémentaires,
que j'utilise via la pile open-source \textbf{Vanetza} (application \code{socktap}) :

\begin{center}
\begin{tabularx}{\textwidth}{@{}l X X@{}}
\toprule
 & \textbf{CAM} & \textbf{DENM} \\
\midrule
Rôle & « je suis là » & « attention, danger ici » \\
Quand & en continu ($\sim$1 Hz) & sur événement \\
Contient & position, cap, vitesse, type & cause + position de l'événement \\
Port BTP & 2001 & 2002 \\
\bottomrule
\end{tabularx}
\end{center}

Le tout fonctionne sur un \textbf{WiFi ordinaire} (diffusion multicast), sans matériel radio
spécialisé (pas de 802.11p).

% =====================================================================
\section{Le CAM --- se signaler en continu}
Le CAM est émis en boucle par chaque véhicule pour annoncer son état.

\subsection{Ce que j'ai amélioré}
Au départ, le CAM portait une position, mais le \textbf{cap} et la \textbf{vitesse} étaient
figés à zéro. Je les ai rendus réels, calculés depuis l'\textbf{odométrie} du robot :
\begin{lstlisting}[language=bash]
odometrie -> cap (lacet) + vitesse (twist)
   -> odom_to_gps.py -> /v2v/my_gps = [lat, lon, cap, vitesse]
   -> v2v_node.py -> socktap -> CAM (Heading, Speed remplis)
\end{lstlisting}

\subsection{Contenu d'un CAM (décodé)}
\begin{center}
\begin{tabularx}{\textwidth}{@{}l l X@{}}
\toprule
\textbf{Champ} & \textbf{Exemple} & \textbf{Signification} \\
\midrule
Station ID & 1 & identifiant du véhicule émetteur. \\
Station Type & 5 & voiture. \\
Latitude / Longitude & $10^{-7}$° & position GPS. \\
Heading & 900 & cap en $0{,}1^\circ$ $\rightarrow$ $90^\circ$. \\
Speed & 500 & vitesse en $0{,}01$~m/s $\rightarrow$ $5$~m/s. \\
\bottomrule
\end{tabularx}
\end{center}

\begin{note}[Vérifié]
Échange CAM \textbf{bidirectionnel} PC$\leftrightarrow$Jetson sur WiFi, avec cap et vitesse
réels issus de l'odométrie. Un moniteur terrain (\code{v2v\_car\_monitor.py}) affiche en
direct la position, le cap, la vitesse et la distance de l'autre véhicule.
\end{note}

% =====================================================================
\section{La perception --- détecter les obstacles}
Pour qu'un véhicule \emph{sache} qu'il y a un danger, il faut le percevoir. J'utilise la
caméra RGB-D (Orbbec Femto Bolt) et un détecteur \textbf{YOLO}.

\subsection{Principe}
\begin{itemize}[leftmargin=1.6em]
  \item \textbf{YOLO} reconnaît les objets (personne, chaise\ldots) sur l'image couleur.
  \item La \textbf{profondeur} donne la distance de chaque objet. Pour un obstacle, je
        prends son point le \textbf{plus proche} (percentile bas), robuste à l'arrière-plan.
  \item Si un objet est trop proche, je lève un signal de \textbf{freinage}
        (\code{/obstacle/brake}).
\end{itemize}

\subsection{Points techniques réglés}
\begin{itemize}[leftmargin=1.6em]
  \item \textbf{Alignement profondeur/couleur} (\code{depth\_registration:=true}) : sans lui,
        les distances étaient fausses (les deux capteurs ont un champ différent).
  \item \textbf{Freinage sur tout obstacle proche}, quelle que soit la classe ; loin = pas de
        freinage. Code couleur à l'écran : rouge = proche, vert = loin.
\end{itemize}

\begin{note}[Vérifié]
Détection en direct sur la caméra de la voiture : une personne à $0{,}4$~m et une chaise
à $0{,}5$~m sont marquées en rouge et déclenchent le freinage ; un mannequin à $4{,}3$~m
reste en vert.
\end{note}

% =====================================================================
\section{Le DENM --- alerter d'un danger}
Quand la perception détecte un obstacle proche, le véhicule émet un \textbf{DENM} pour
prévenir les autres.

\subsection{Ce que j'ai ajouté}
\code{socktap} ne savait émettre que des CAM. J'ai ajouté une application DENM
(\code{denm\_application.cpp}) qui construit un vrai DENM ETSI et l'émet sur commande :
\begin{itemize}[leftmargin=1.6em]
  \item \textbf{cause} : $99$ = situation dangereuse, sous-cause $2$ = freinage d'urgence ;
  \item \textbf{position} de l'événement (celle du véhicule) ;
  \item \textbf{déclenchement} par un message UDP, envoyé par ROS quand
        \code{/obstacle/brake} passe à vrai.
\end{itemize}
À la réception, le DENM est décodé et publié sur \code{/v2v/denm}, puis affiché comme une
\textbf{alerte}.

\subsection{CAM ou DENM ?}
\begin{center}
\begin{tabularx}{\textwidth}{@{}l X@{}}
\toprule
\textbf{Le CAM} & répond à « qui est autour de moi, où, à quelle vitesse ? » (permanent) \\
\textbf{Le DENM} & répond à « que se passe-t-il d'anormal, où, maintenant ? » (ponctuel) \\
\bottomrule
\end{tabularx}
\end{center}

\begin{note}[Vérifié]
DENM émis (déclenché par le freinage) et \textbf{reçu} sur l'autre véhicule : cause $99/2$,
position et distance affichées.
\end{note}

% =====================================================================
\section{La démonstration de bout en bout}
J'ai relié toutes les briques : la voiture perçoit un obstacle, émet un DENM, et l'autre
véhicule (le PC) reçoit et affiche l'alerte.

\subsection{La chaîne complète}
\begin{lstlisting}
Camera (Jetson) --> YOLO + profondeur (PC) --> /obstacle/brake
                                                    |
                                                    v
                      v2v_node "voiture A" --> DENM  (Vanetza / WiFi)
                                                    |
                                                    v
              v2v_node "voiture B" (PC) --> /v2v/denm --> *** ALERTE ***
\end{lstlisting}

\subsection{Comment la lancer}
Deux domaines ROS séparés, pour que l'alerte passe bien par le V2V (Vanetza) et non par ROS.

\begin{lstlisting}[language=bash]
# Jetson (voiture A - domaine 125)
./launch_camera.sh        # camera (profondeur alignee)
./launch_denm_car.sh      # emet un DENM quand un obstacle est detecte

# PC
./run_yolo_pc.sh          # (domaine 125) detection YOLO -> /obstacle/brake
./launch_denm_pc.sh       # (domaine 2)   recoit le DENM + AFFICHE l'alerte
\end{lstlisting}
Un obstacle devant la caméra $\rightarrow$ le PC affiche :
\begin{lstlisting}
*** ALERTE V2V ***   DENM recu du vehicule 2
Situation dangereuse (freinage d'urgence)   Distance : 13.3 m
\end{lstlisting}

% =====================================================================
\section{Bilan et perspectives}
\textbf{Ce qui marche} : la boucle de sécurité coopérative complète --- percevoir (YOLO +
profondeur), décider (freinage), alerter (DENM ETSI), recevoir et afficher --- sur WiFi,
entre deux véhicules, avec des messages CAM et DENM conformes.

\textbf{Perspectives} :
\begin{enumerate}[leftmargin=1.6em]
  \item \textbf{Embarquer la perception} : porter YOLO sur la Jetson pour que la voiture
        détecte et alerte toute seule (freinage d'urgence réel à bord).
  \item \textbf{Réaction physique} : buzzer / LED, puis freinage automatique (AEB) via l'ESC.
  \item \textbf{Démo sur 5G} : rejouer l'échange (le cellulaire bloque le multicast, il
        faudra passer en unicast).
  \item \textbf{Sécurité} : messages \textbf{signés} (\code{--security certs}, vrai C-ITS).
\end{enumerate}

\vspace{6pt}\hrule
{\footnotesize\itshape Fichiers --- CAM/V2V : \code{v2v\_node.py}, \code{odom\_to\_gps.py},
\code{v2v\_car\_monitor.py} ; perception : \code{yolo\_detect.py}, \code{run\_yolo\_pc.sh} ;
DENM : \code{denm\_application.cpp/.hpp} (socktap), \code{denm\_alert.py},
\code{launch\_denm\_car.sh}, \code{launch\_denm\_pc.sh}. Docs liés :
\code{rapport\_complet\_v2v\_cam.tex}, \code{v2v\_explique\_pas\_a\_pas.tex},
\code{demo\_terrain\_v2v.tex}.}

\end{document}
